% Chapter 13, generated by make_chapters.py from papers/Ch13_Bio_09_Polyploid_Phasing/anccft22.tex.
% Source paper: Bio 9
\chapter[Polyploid phasing]{\texorpdfstring{Polyploid phasing as a $p$-fold decomposition of a balanced flow, and the homogeneous space of alleles}{Polyploid phasing as a p-fold decomposition of a balanced flow, and the homogeneous space of alleles}}\label{chap:ch13}
\begin{chaptersummary}
The $k$-mer spectrum of $p$ homologous chromosomes sequenced together is the arc
multiplicity vector $m=\sum_i\spec(H_i)$ of one pooled de Bruijn multigraph, and phasing is
the recovery of the summands. We show that the labelled phasings --- ordered $p$-tuples of
circular sequences whose spectra add to $m$ --- are counted by a sum over the ordered
decompositions $m=c_1+\dots+c_p$ into connected balanced points of the box $[0,m]$ of the
circulation lattice, each term the product of the per-sequence BEST counts
$\Nseq(c_i)=t(c_i)\prod(d^+-1)!/\prod c_i(a)!$. This is the lattice sum of the preceding
note with the strand involution removed and one constraint replaced by a $p$-fold partition.
The symmetric group $\Sym_p$ permutes the parts of a decomposition; it does \emph{not} act on
the pooled graph, and the natural map from the disjoint union of the haplotype graphs to the
pooled graph is a covering only when the haplotypes share no $k$-mer, so there is no
monodromy in the sense of covering spaces. What replaces it is a homogeneous space: at a
biallelic site carried by $a$ of the $p$ haplotypes the local choice is an $a$-subset of the
labels, $\Sym_p/(\Sym_a\times\Sym_{p-a})$, and for $b$ sites in distinct $(k{+}1)$-windows on a
repeat-free reference the labelled count is $\prod_j\binom{p}{a_j}$ exactly, which
$|\Hom(\pi_1,\Sym_p)|=(p!)^b$ does not reproduce. Sites within one $(k{+}1)$-window merge into
a single multi-allelic site whose alleles are the distinct local haplotypes. Unlabelled
phasings are multisets, counted by $\prod_c\binom{\Nseq(c)+r_c-1}{r_c}$ over the multiset of
spectra, equal to the labelled count divided by $p!$ exactly when the parts are pairwise
distinct. Dropping the requirement that all haplotypes have the reference length exposes the
copy-number ambiguity through repeats: the pooled spectrum of two copies of $\phi$X174 is also
the pooled spectrum of two circles of $2496$ and $8276$ base pairs, and of two tandem-dimer
circles, whose weight in the sum is $1/2$. Every count is verified on $\phi$X174 with synthetic
substitutions at $p=2$ and $p=3$ by an enumeration of walks that does not use the lattice.
\end{chaptersummary}

\medskip
\noindent\textbf{Keywords.} de Bruijn graph, haplotype phasing, polyploid, BEST theorem,
circulation lattice, flow decomposition, homogeneous space, copy-number variation.

\noindent\textbf{MSC 2020.} 05C20, 05C21, 05C45, 05C50, 92D20.

\section{What was asked, and what is true}

The preceding note Chapter~\ref{chap:ch07} showed that a double-stranded reconstruction of a circular
sequence $S$ is a closed walk through half the arcs of the reverse-complement double cover
whose image under the strand involution $\rho$ walks the other half, and that the count is a
sum over the box points of the $(-1)$-eigenlattice of $\rho$ on circulations. The involution
was the obstacle: it reverses arcs, it is a duality rather than a symmetry of the digraph, and
all the machinery of covering graphs is unavailable.

Polyploid phasing looks at first like the case where that obstacle is absent. The pooled
$k$-mer spectrum of $p$ homologous circular chromosomes $H_1,\dots,H_p$ is the arc
multiplicity vector $m=\sum_i\spec(H_i)$ of the pooled de Bruijn multigraph, and a phasing is
a way of writing $m$ as a sum of $p$ spectra of circular sequences. Exchanging two haplotypes
is a permutation of the summands; it does not reverse anything. One is tempted to say that the
symmetric group $\Sym_p$ acts on the pooled graph as a genuine automorphism group, that the
haplotypes are the sheets of a $p$-fold cover, and that the phasing space is the space of flat
$\Sym_p$-connections $H^1(G_k,\Sym_p)=\Hom(\pi_1(G_k),\Sym_p)$ on the pooled graph, with a
transposition at each heterozygous site. The specification of this note said so, and so did
its planning.

None of that is so, and the correct statements are simpler. $\Sym_p$ acts on the set of
\emph{decompositions} of $m$, not on the graph: the pooled graph has no distinguished symmetry
exchanging haplotypes, because the haplotypes are different words and the graph does not
remember which arc came from which. The natural map from the disjoint union of the haplotype
graphs to the pooled graph is a covering map only when no $k$-mer is shared, in which case
the pool is the disjoint union and there is nothing to phase (Proposition~\ref{ch13:prop:cover}).
And the local structure at a biallelic site is not a group but a homogeneous space: if $a$ of
the $p$ haplotypes carry one allele and $p-a$ the other, the choice at the site is which $a$
labels carry it, a point of $\Sym_p/(\Sym_a\times\Sym_{p-a})$ of cardinality $\binom pa$; only
for $p=2$, or at a site where all $p$ alleles differ, is this the group. For $b$ sites in
distinct $(k{+}1)$-windows on a repeat-free reference the labelled count is exactly
$\prod_j\binom{p}{a_j}$ (Theorem~\ref{ch13:thm:sites}), and $(p!)^b$ is wrong already at $p=3$,
$b=1$, $a=1$: the answer is $3$, not $6$.

What survives, and is the subject of Section~\ref{ch13:sec:lattice}, is the lattice sum. The
labelled phasings are counted exactly by
\[
  \Lab(H_1{+}\dots{+}H_p,k)=\sum_{\substack{c_1+\dots+c_p=m\\ c_i\in Z_1\cap[0,m],\ \supp c_i
  \text{ connected}}}\ \prod_{i=1}^p\Nseq(c_i),
\]
with $\Nseq$ the per-sequence BEST count of Chapter~\ref{chap:ch07}; the unlabelled phasings by a multiset
formula over the unordered decompositions (Section~\ref{ch13:sec:unlabelled}); and when the parts
are not required to have the reference length the sum exposes the copy-number ambiguity
through repeats, including tandem-dimer circles whose weight is $1/2$ (Section~\ref{ch13:sec:cnv}).
Section~\ref{ch13:sec:data} records the verification on $\phi$X174.

Throughout, sequences are circular over $\{A,C,G,T\}$, a \emph{circular sequence} is a word
up to rotation, and $d(T)$ is the number of rotations fixing $T$. All computations are exact
integer or rational arithmetic; the script \texttt{polyploid\_phasing.py} reproduces every
number below and prints a six-line pass/fail battery.

\section{The pooled multigraph and the phasing lattice}\label{ch13:sec:lattice}

\begin{definition}\label{ch13:def:pool}
Let $H_1,\dots,H_p$ be circular sequences of lengths $N_1,\dots,N_p$. The \emph{pooled de
Bruijn multigraph} $G_k=G_k(H_1{+}\dots{+}H_p)$ has one vertex per $k$-mer occurring in some
$H_i$ and one arc per occurrence of a $(k{+}1)$-mer in some $H_i$. A vertex is a \emph{branch
vertex} if it has at least two distinct successors or at least two distinct predecessors in
the pool; the \emph{compacted} pooled graph has the branch vertices as vertices and the
maximal unitigs between them as arcs. An arc is labelled by its haplotype and position; its
\emph{content} is the unitig string, $A$ is the set of contents, $m(a)$ the number of arcs of
content $a$, $\ell(a)$ its number of $(k{+}1)$-mers, and $\partial\colon\Z^A\to\Z^V$ the
boundary. $Z_1=\ker\partial$ is the circulation lattice.
\end{definition}

The compaction criterion differs from the one used for a single genome in Chapter~\ref{chap:ch15}, where a
$k$-mer is a branch vertex if and only if it is repeated. In a pool every shared $k$-mer is
repeated, so that criterion would contract nothing; the criterion of
Definition~\ref{ch13:def:pool} contracts exactly the $k$-mers through which every haplotype must
pass in the same way. The per-sequence count below is intrinsic and does not depend on which
compaction is used.

\begin{definition}\label{ch13:def:nseq}
For $c\in\Z^A_{\ge0}$ balanced with connected support let $G_c$ be the multigraph with $c(a)$
labelled parallel arcs of content $a$, $t(c)$ its number of spanning arborescences
converging on a root, $\Loc(c)=\prod_v(d^+_c(v)-1)!$ over the vertices of positive $c$-degree,
and
\[
  \Nseq(c)=\frac{t(c)\,\Loc(c)}{\prod_a c(a)!}\;=\;\sum_{\spec T=c}\frac1{d(T)}\in\Q_{>0},
\]
the second equality being Chapter~\ref{chap:ch07}, Thm.~6. $\Nseq(c)=0$ if the support is disconnected.
The \emph{length} of $c$ is $|c|=\sum_a c(a)\ell(a)$.
\end{definition}

\begin{lemma}\label{ch13:lem:spectrum}
Let $T$ be a circular sequence all of whose $(k{+}1)$-mers occur in the pool. Then $T$
traverses each unitig of the compacted pooled graph a well-defined number of times, so
$\spec(T)$ is a vector $c\in\Z^A_{\ge0}$ with $\partial c=0$ and connected support, and
$|c|$ is the length of $T$.
\end{lemma}

\begin{proof}
An interior $k$-mer $w$ of a unitig has exactly one successor and one predecessor in the pool.
Every $(k{+}1)$-mer of $T$ through $w$ is a $(k{+}1)$-mer of the pool, hence is the unique
in- or out-arc of $w$; so each passage of $T$ through $w$ continues along the unitig in both
directions, and the number of passages is the same at every $k$-mer of the unitig.
\end{proof}

\begin{definition}\label{ch13:def:phasing}
A \emph{labelled phasing} of the pool at level $k$ is an ordered $p$-tuple $(T_1,\dots,T_p)$
of circular sequences with $\sum_i\spec(T_i)=m$ as multisets of $(k{+}1)$-mers; its
\emph{weight} is $\prod_i 1/d(T_i)$. An \emph{unlabelled phasing} is the multiset
$\{T_1,\dots,T_p\}$. An \emph{ordered decomposition} of $m$ is a $p$-tuple
$(c_1,\dots,c_p)$ of balanced vectors with connected support, $0\le c_i\le m$ and
$\sum c_i=m$; $\mathcal D_p(m)$ is their set. The \emph{equal-length} phasings and
decompositions are those with $|T_i|=N_1$, resp.\ $|c_i|=N_1$, for all $i$.
\end{definition}

\begin{theorem}[The phasing lattice]\label{ch13:thm:lattice}
The map $(T_1,\dots,T_p)\mapsto(\spec T_1,\dots,\spec T_p)$ is a surjection from the labelled
phasings onto $\mathcal D_p(m)$, and the weighted number of labelled phasings is
\[
  \Lab(H_1{+}\dots{+}H_p,k)\;=\;\sum_{(c_1,\dots,c_p)\in\mathcal D_p(m)}\ \prod_{i=1}^p\Nseq(c_i).
\]
The same holds for equal-length phasings with the sum restricted to equal-length
decompositions. In particular the count is an integer whenever every part of every
decomposition is the spectrum of aperiodic sequences only.
\end{theorem}

\begin{proof}
Lemma~\ref{ch13:lem:spectrum} sends each $T_i$ to a balanced connected $c_i\ge0$, and
$\sum\spec T_i=m$ is $\sum c_i=m$, which also gives $c_i\le m$. Conversely each part of a
decomposition is balanced and connected, so $G_{c_i}$ has an Eulerian circuit, which read as
a sequence has spectrum $c_i$; the $p$ sequences together have spectrum $m$. The fibre over
$(c_1,\dots,c_p)$ is the product of the fibres over the $c_i$, each of weighted cardinality
$\Nseq(c_i)$ by Definition~\ref{ch13:def:nseq}. Lengths are preserved by Lemma~\ref{ch13:lem:spectrum}.
\end{proof}

\begin{remark}
The lattice is the same object as in Chapter~\ref{chap:ch07}: there $p=2$, $c_2=\rho c_1$ is forced, and
the second part is the strand image of the first, so $\mathcal D_2(m)$ collapses to the box
points of the anti-invariant lattice. Here nothing ties the parts together except $\sum c_i=m$,
and $\Sym_p$ acts on $\mathcal D_p(m)$ by permuting them. The single-strand theory is
$p=1$, where $\mathcal D_1(m)=\{m\}$ and the sum is one term.
\end{remark}

\begin{proposition}[No covering]\label{ch13:prop:cover}
Let $\pi\colon\bigsqcup_i G_k(H_i)\to G_k(H_1{+}\dots{+}H_p)$ send each vertex and arc of a
haplotype graph to the same $k$-mer or occurrence in the pool. Then $\pi$ is a surjective
homomorphism of multidigraphs, and it is a covering map if and only if no $k$-mer occurs in
two different haplotypes, in which case $\pi$ is an isomorphism. In particular, when
$p\ge2$ and the haplotypes share a $k$-mer, $\Sym_p$ acts on $\mathcal D_p(m)$ but there is no
$p$-sheeted cover of the pooled graph with deck group $\Sym_p$, and no action of $\Sym_p$ on
the pooled graph is induced by relabelling haplotypes.
\end{proposition}

\begin{proof}
A covering map is a local isomorphism: the out-arcs of a vertex $\tilde w$ upstairs map
bijectively onto the out-arcs of $\pi(\tilde w)$. If a $k$-mer $w$ occurs in $H_i$ and in
$H_j$, $i\ne j$, then $w$ has two preimages $\tilde w_i,\tilde w_j$, the out-degree of
$\tilde w_i$ is the multiplicity of $w$ in $H_i$, and the out-degree of $w$ is the sum of the
multiplicities of $w$ in all haplotypes, which strictly exceeds it. If no $k$-mer is shared,
$\pi$ is a bijection on vertices and on arcs. For the last sentence: an action of $\Sym_p$ on
the pooled graph induced by relabelling would have to fix every vertex and every arc, since
the pool does not record labels, and the only such action is trivial.
\end{proof}

The honest picture is therefore this. The pooled graph carries the flow $m$; a phasing is a
decomposition of the flow into $p$ closed walks; the labels are a decoration of the
decomposition and $\Sym_p$ permutes the decoration. The question ``how many phasings'' is a
question about the fibres of $\mathcal D_p(m)\to\mathcal D_p(m)/\Sym_p$, and it is answered
locally by homogeneous spaces (Section~\ref{ch13:sec:sites}) and globally by a multiset count
(Section~\ref{ch13:sec:unlabelled}).

\section{The homogeneous fibre, and linkage}\label{ch13:sec:sites}

Fix a circular \emph{reference} $R$ of length $N$ and let each $H_i$ be obtained from $R$ by
substitutions at a set of positions, the \emph{sites}; at a site $s$ let the alleles be the
letters occurring there across the $H_i$ and $a_s(\alpha)$ the number of haplotypes carrying
allele $\alpha$. A site is \emph{biallelic} if two letters occur, and we then write $a_s$ for
the count of the minor allele. Two sites are in the same \emph{$(k{+}1)$-window} if some
$(k{+}1)$-mer position of $R$ contains both, i.e.\ their circular distance is at most $k$.

\begin{theorem}[Independent sites on a repeat-free reference]\label{ch13:thm:sites}
Suppose $R$ has no repeated $k$-mer, the sites lie in pairwise distinct $(k{+}1)$-windows, and
no $k$-mer of any $H_i$ overlapping a site coincides with a $k$-mer of $R$ or with one
overlapping another site. Then the equal-length ordered decompositions of $m$ are in bijection
with $\prod_s\Sym_p/\big(\prod_\alpha\Sym_{a_s(\alpha)}\big)$, every part has $\Nseq=1$, and
\[
  \Lab=\prod_s\binom{p}{a_s(\alpha_1),\,a_s(\alpha_2),\,\dots}
  \;=\;\prod_{s\ \text{biallelic}}\binom{p}{a_s}\quad\text{when all sites are biallelic.}
\]
\end{theorem}

\begin{proof}
Under the hypotheses the compacted pooled graph is a cycle of $N$ base pairs in which each
site $s$ is replaced by a \emph{bubble}: the $k$-mer of $R$ ending just before $s$ is a branch
vertex $u_s$, the $k$-mer starting just after $s$ is a branch vertex $v_s$, and between them
run parallel unitigs, one per allele $\alpha$ at $s$, each of length $k{+}1$ in
$(k{+}1)$-mers and of multiplicity $a_s(\alpha)$; the $k$-mers strictly between sites are
shared by all haplotypes and contract into unitigs of multiplicity $p$. There are no other
branch vertices because $R$ has no repeats and the allele $k$-mers are new. A closed walk in
this graph has a winding number $w\ge1$ around the cycle and length exactly $wN$, since all
branches of a bubble have the same length. A part of an equal-length decomposition has length
$N$, hence $w=1$, hence passes each bubble exactly once through one branch; its spectrum is
the $0/1$ vector selecting one allele per site, $G_c$ is a simple cycle, and $\Nseq(c)=1$.
An ordered decomposition is thus a $p$-tuple of allele selections, one per label, with
$a_s(\alpha)$ labels selecting $\alpha$ at each site: exactly the product over sites of the
sets of ordered partitions of the labels into classes of sizes $a_s(\alpha)$, which are the
homogeneous spaces stated, of cardinalities the multinomial coefficients. Theorem~\ref{ch13:thm:lattice}
gives the count.
\end{proof}

\begin{corollary}\label{ch13:cor:notH1}
Under the hypotheses of Theorem~\ref{ch13:thm:sites} the pooled graph has $\pi_1$ free of rank
$\sum_s(n_s-1)$, where $n_s$ is the number of alleles at $s$, and
$|\Hom(\pi_1,\Sym_p)|=(p!)^{\rank\pi_1}$, which equals $\Lab$ only when $p\le2$ or every
site has $p$ distinct alleles.
\end{corollary}

\begin{proof}
The cycle rank of a cycle with bubbles is the number of extra branches. For $p=2$ and biallelic
sites both counts are $2^b$; for $p\ge3$ and a biallelic site with $a_s=1$ the site
contributes $p$ to $\Lab$ and $p!$ to the other.
\end{proof}

\begin{theorem}[Linkage]\label{ch13:thm:linkage}
Let a set $B$ of sites lie in one $(k{+}1)$-window, with the hypotheses of
Theorem~\ref{ch13:thm:sites} otherwise in force for $B$ as a whole (no $k$-mer overlapping $B$ is a
$k$-mer of $R$ or overlaps another site). Then in the compacted pooled graph $B$ is a single
bubble whose branches are the distinct \emph{local haplotypes} --- the distinct restrictions
of $H_1,\dots,H_p$ to $B$ --- with multiplicity the number of haplotypes carrying each, and in
Theorem~\ref{ch13:thm:sites} the block $B$ counts as one site whose alleles are its local haplotypes.
In particular two biallelic sites within distance $k$ carried in only two combinations
contribute $\binom{p}{a}$ and not $\binom{p}{a}\binom{p}{a'}$.
\end{theorem}

\begin{proof}
Some $(k{+}1)$-mer position of $R$ contains all of $B$, so every haplotype has a
$(k{+}1)$-mer through the whole block, and two haplotypes with the same local restriction have
the same $k$-mers throughout the window while two with different restrictions differ in that
$(k{+}1)$-mer. The $k$-mers overlapping $B$ but not containing all of it may be shared between
haplotypes that agree on the sites they contain; this merges vertices inside the bubble but
creates no new closed walk of length $N$, because a walk of length $N$ has winding number one
and its passage through the bubble is determined by the $(k{+}1)$-mer it uses at the position
containing all of $B$. So the equal-length parts through the block are the local haplotypes,
and the count is the multinomial in their multiplicities.
\end{proof}

For sites that are pairwise within distance $k$ but not all inside one $(k{+}1)$-window, or a
chain of sites each within $k$ of the next, the bubble structure is a chain of partially merged
bubbles and we do not have a closed form; see the open questions.

\begin{remark}[With repeats]\label{ch13:rem:repeats}
When $R$ has repeats the argument of Theorem~\ref{ch13:thm:sites} gives a lower bound. Each allele
selection $c$ --- $\spec(R)$ with one allele chosen at each site --- is still a part of length
$N$, and $G_c$ is isomorphic to $G_{\spec R}$ with the bubbles replaced by single unitigs of
the same length, so $\Nseq(c)=\Nseq(R)$; the ordered allele selections therefore contribute
$\prod_s\binom{p}{a_s}\,\Nseq(R)^p$ to $\Lab$. Other equal-length decompositions can exist only
if $p\,\spec(R)$ admits a balanced connected sub-vector of length $N$ that is not an allele
selection, which requires a coincidence among the lengths of the sub-circles of $R$ through
its repeats. For $\phi$X174 at $k=12$ no such coincidence occurs and the bound is the count at
every tested configuration (Table~\ref{ch13:tab:diploid}).
\end{remark}

\section{The unlabelled multiset formula and the
\texorpdfstring{$\Sym_p$}{Sp}-action}\label{ch13:sec:unlabelled}

\begin{proposition}\label{ch13:prop:multiset}
$\Sym_p$ acts on $\mathcal D_p(m)$ by permuting the parts; the stabiliser of
$(c_1,\dots,c_p)$ is $\prod_c\Sym_{r_c}$, where $r_c$ is the number of parts equal to $c$.
If every sequence with spectrum a part of the decomposition is aperiodic, the number of
unlabelled phasings whose spectra form the multiset $\{c^{r_c}\}$ is
\[
  \prod_{c}\binom{\Nseq(c)+r_c-1}{r_c},
\]
and the total number of unlabelled phasings is the sum of this over the orbits
$\mathcal D_p(m)/\Sym_p$. It equals $\Lab/p!$ exactly when every decomposition has pairwise
distinct parts.
\end{proposition}

\begin{proof}
An unlabelled phasing with spectrum multiset $\{c^{r_c}\}$ is, for each $c$, a multiset of
$r_c$ sequences drawn from the $\Nseq(c)$ sequences with spectrum $c$ (integers, by
aperiodicity); multisets of size $r$ from a set of size $n$ number $\binom{n+r-1}{r}$. The
labelled phasings over an orbit number $|{\rm orbit}|\cdot\prod_c\Nseq(c)^{r_c}=
(p!/\prod_c r_c!)\prod_c\Nseq(c)^{r_c}$, and $\binom{n+r-1}{r}=n^r/r!$ only for $r\le1$.
\end{proof}

The distinction is not academic. For three haplotypes $S,S',S$ differing at one site, the
three ordered decompositions all have the same multiset of spectra, and the unlabelled count is
$\Nseq(S)\binom{\Nseq(S)+1}{2}$, which is $6$ for $\phi$X174 at $k=12$, against a labelled
count of $24$; $24/6=4\ne6$.

\begin{remark}[Periodic parts]\label{ch13:rem:periodic}
When a part $c$ is $2c'$ for a circulation $c'$ and $G_{c'}$ is Eulerian, the sequences with
spectrum $c$ include the tandem dimers $UU$ of the sequences $U$ with spectrum $c'$, which
have $d(UU)=2$, and more generally every sequence with spectrum $c$ may be periodic, in which
case $\Nseq(c)$ is a half-integer or worse. The weighted count of Theorem~\ref{ch13:thm:lattice}
remains correct, the multiset formula does not apply to such parts, and the script reports
them separately. In the data they arise only when the equal-length constraint is dropped, and
they are the dimer circles of Section~\ref{ch13:sec:cnv}.
\end{remark}

\section{Copy-number variation across repeats}\label{ch13:sec:cnv}

The equal-length constraint encodes the substitution-only model. Without it, the sum of
Theorem~\ref{ch13:thm:lattice} ranges over all decompositions of the pooled flow into closed walks,
and for a reference with repeats it contains decompositions in which a block bounded by two
copies of a repeat is carried once by one haplotype and three times by the other, or in which
the two haplotypes are two circles of different sizes. These are the copy-number and
excision--integration ambiguities of the pooled spectrum, and the lattice sum enumerates them.

\begin{sloppypar}
For $\phi$X174 at $k=12$ the single-strand graph is the interleaved two-copy repeat of
Chapter~\ref{chap:ch01}, $a\,X_1\,b\,Y_1\,a\,X_2\,b\,Y_2$: two branch $12$-mers $a,b$ and four unitigs,
with $\Nseq(S)=2$. The sub-circles $aX_ibY_j$ have lengths $2496$, $2890$, $2638$ and $2748$,
in two complementary pairs summing to $N=5386$. The pool $S+S$ has $m=2$ on each unitig, and
its ordered decompositions with connected parts number $17$: one with two parts of length $N$,
and sixteen of the shapes
\end{sloppypar}
\[
  (2496,\,8276),\ (2890,\,7882),\ (2638,\,8134),\ (2748,\,8024),\ (4992,\,5780),\ (5276,\,5496)
\]
and their transposes. The first four are a sub-circle $aX_ibY_j$ of length $\ell$ and the complementary circle of
length $2N-\ell$, which carries the two unitigs of the sub-circle once and the other two
twice: a sub-circle excised from one copy of the genome and its content integrated into the
other. The last two
have \emph{both} parts periodic: $4992=2\cdot2496$ and $5780=2\cdot2890$, the tandem dimers of
a complementary pair of sub-circles, each with $d(T)=2$ and weight $1/2$. With one substitution
added the decompositions number $24$, of which $22$ have unequal lengths
(Table~\ref{ch13:tab:cnv}). A pooled spectrum thus does not distinguish two identical circular
chromosomes from two dimer circles of complementary halves, and the count records the fact
with the weight it deserves.

\section{Verification and data}\label{ch13:sec:data}

All configurations use $\phi$X174 (NC\_001422.1, $5386$\,bp) as the reference at $k=12$, where
$\Nseq(S)=2$, with substitutions at positions $600,1700,2900,3900,4900$ (pairwise distance
$>k$, none within $k$ of a repeat $k$-mer) and, for linkage, at $606$. Each row of
Tables~\ref{ch13:tab:diploid}--\ref{ch13:tab:cnv} is computed twice: from the lattice, by enumerating the
box $Z_1\cap[0,m]$ with the balance constraint and evaluating $\Nseq$ at every point; and by an
enumeration of closed walks $T_1$ in the content multigraph (rooted at the smallest content
$a_0$, weighted $1/u(a_0)$, usage bounded by $m$) that never sees the lattice, followed for
$p=2$ by an explicit enumeration of the partner walks $T_2$ and a count of the unordered pairs
$\{T_1,T_2\}$ as sets.

\begin{table}[t]
\centering
\small
\begin{tabular}{lrrrrrrrr}
\toprule
configuration & $|V|$ & $|A|$ & $\rank Z_1$ & box & conn. & decomp. & $\Lab$ & $\Unl$\\
\midrule
1 site              &  4 &  7 & 4 &  26 &  25 &  2 &   8 &  4\\
2 sites             &  6 & 10 & 5 &  40 &  39 &  4 &  16 &  8\\
3 sites             &  8 & 13 & 6 &  56 &  55 &  8 &  32 & 16\\
5 sites             & 12 & 19 & 8 & 122 & 121 & 32 & 128 & 64\\
$600,606$ linked $+$ 2 sites & 8 & 13 & 6 & 56 & 55 & 8 & 32 & 16\\
$\{600,606\}$ vs $\{600\}$   & 4 &  7 & 4 & 26 & 25 & 2 &  8 &  4\\
\bottomrule
\end{tabular}
\caption{Diploid pools, equal length. ``decomp.'' is the number of equal-length ordered
decompositions; the prediction is $2^b$ for $b$ linkage blocks, $\Lab=2^b\Nseq(S)^2$ and
$\Unl=2^{b-1}\Nseq(S)^2$, met in every row. The walk enumeration reproduces $\Nseq$ on every
spectrum it finds and the same $\Lab$; the explicit unordered pairs number $4,8,16,64,16,4$.
Checks (P), (L), (W), (U).}
\label{ch13:tab:diploid}
\end{table}

\begin{table}[t]
\centering
\small
\begin{tabular}{lrrrrrrrr}
\toprule
configuration & $|V|$ & $|A|$ & $\rank Z_1$ & box & conn. & decomp. & $\Lab$ & $\Unl$\\
\midrule
three sites: $a=1$, $a=2$, triallelic & 8 & 14 & 7 & 248 & 247 & 54 & 432 & 72\\
one site, haplotypes $S,S',S$ & 4 & 7 & 4 & 68 & 67 & 3 & 24 & 6\\
\bottomrule
\end{tabular}
\caption{Triploid pools, equal length. First row: sites $600$ (minor allele in one haplotype),
$1700$ (in two) and $2900$ (three alleles); $\binom31\binom32\,3!=54$ decompositions, all
with pairwise distinct parts, $\Lab=54\cdot2^3$ and $\Unl=\Lab/6$; $|\Hom(\pi_1,\Sym_3)|$ would
be $6^3=216$. Second row: three decompositions with the same multiset of spectra,
$\Unl=\Nseq(S)\binom{\Nseq(S)+1}{2}=6\ne24/6$. Check (T).}
\label{ch13:tab:triploid}
\end{table}

\begin{table}[t]
\centering
\small
\begin{tabular}{lrrrrrr}
\toprule
configuration & box & conn. & decomp. & unequal & periodic & $\Lab$\\
\midrule
diploid, 1 site, all lengths & 26 & 25 & 24 & 22 & 4 & 66\\
homozygous $S+S$, all lengths & 19 & 18 & 17 & 16 & 4 & 33\\
\bottomrule
\end{tabular}
\caption{Without the equal-length constraint. ``unequal'' is the number of ordered
decompositions whose parts have different lengths; ``periodic'' the number with a periodic
part, each counted with weight $1/2$ per such part. Length profiles for $S+S$: one $(N,N)$,
and $(2496,8276)$, $(2890,7882)$, $(2638,8134)$, $(2748,8024)$, $(4992,5780)$, $(5276,5496)$
twice each. The explicit pair enumeration for the diploid row finds $34$ unordered pairs: $32$
on the ten decompositions without a periodic part, matching the multiset formula, and two
pairs of dimer circles. Check (C).}
\label{ch13:tab:cnv}
\end{table}

\paragraph{The six checks.} (P) the pooled compacted graph is balanced and gains exactly two
branch vertices per site away from repeats; (L) $b$ unlinked biallelic sites give $2^b$
equal-length decompositions and $2^b\Nseq(S)^2$ labelled phasings, and two sites within $k$
are one site; (W) for $p=2$ the walk enumeration reproduces $\Nseq$ on every spectrum and
$\sum_{T_1}\Nseq(m-\spec T_1)$ equals the lattice sum; (U) the unordered phasings enumerated
as sets of explicit walk pairs equal the multiset formula; (T) for $p=3$ the decompositions
number $\prod_s\binom{3}{a_s}$, $3!$ at the triallelic site, and the unlabelled count is the
multiset count; (C) free lengths admit copy-number decompositions through the repeat, and the
equal-length unlabelled count of $S+S$ is $\binom{\Nseq(S)+1}{2}=3$. All six pass, in under a
second.

\section{Scope, provenance, and what is open}

\paragraph{Scope.} Theorem~\ref{ch13:thm:lattice}, Proposition~\ref{ch13:prop:cover} and
Proposition~\ref{ch13:prop:multiset} hold for arbitrary circular haplotypes. Theorem~\ref{ch13:thm:sites},
Corollary~\ref{ch13:cor:notH1} and Theorem~\ref{ch13:thm:linkage} assume a repeat-free reference and
isolated sites; Remark~\ref{ch13:rem:repeats} is a lower bound for references with repeats, met with
equality in all data. The model is the pooled $k$-mer spectrum alone: no reads, no read
pairs, no coverage, no errors. Reads that span sites constrain the decompositions further and
are the whole content of the haplotype assembly problem \cite{LBIL01,BB08,AI12}; what is
counted here is the ambiguity that remains before any read is used, which is the analogue for
phasing of the reconstruction count of \cite{KSP} for assembly.

\paragraph{Provenance.} The BEST theorem is \cite{BEST,SmithTutte}, its per-sequence
normalisation and the notation $\Nseq$ are Chapter~\ref{chap:ch07}, and the single-genome decomposition
of $\phi$X174 at $k=12$ is Chapter~\ref{chap:ch01}. Haplotype assembly as a combinatorial problem on
sites is \cite{LBIL01}; the cycle-basis formulation of \cite{AI12} is the closest in spirit to
the circulation lattice here, though it lives on a graph of sites and reads rather than on the
de Bruijn graph. Polyploid phasing methods are surveyed in \cite{Garg21}. We are not aware of
prior work counting phasings through the de Bruijn graph, or of the homogeneous-space
description of the local fibre; a targeted search found none.

\paragraph{Open.} \emph{First}, the count for a chain of sites each within $k$ of the next
but not all in one window (partially merged bubbles). \emph{Second}, the equality in
Remark~\ref{ch13:rem:repeats}: characterise the references for which $p\,\spec(R)$ has no
balanced connected sub-vector of length $N$ beyond the allele selections. \emph{Third}, the
copy-number decompositions at genome scale, where the lattice has the rank of $Z_1$ of the
pooled skeleton and the sum is not enumerable; a transfer-matrix evaluation along the genome
is the natural tool, as it is for the double-stranded sum of Chapter~\ref{chap:ch07}. \emph{Fourth}, the
effect of reads: a read spanning sites $s,s'$ is a linear constraint on the decomposition,
and the phasings compatible with a read set are a sub-lattice section; whether the standard
phasing objectives (minimum error correction, maximum fragments cut) have an expression in
these terms is unexplored.
